\chapter{Problem Statement}

We consider a distributed RL setup where an agent learns to control the CarRacing-v3 environment through interaction. The agent consists of a convolutional neural network (CNN) for feature extraction and an actor-critic network for policy and value estimation. Given resource constraints, we partition this agent across two machines: Machine~0 hosts the CNN and environment interaction, while Machine~1 hosts the actor-critic network.

\section{Communication Overhead Analysis}

As an example, we compare data transfer requirements across two deployment scenarios over 100 training iterations with batch size 4096 (100 iterations $\times$ 4096 steps per iteration).

\textbf{Cloud Setup:} The local machine transmits raw observations to the cloud for complete training. Each observation is a $4 \times 96 \times 96$ tensor. Total data transfer:
\begin{equation}
100 \times 4096 \times (4 \times 96 \times 96) = 14.7\text{B tensors} \approx 28\text{ GB}
\end{equation}

\textbf{Pipeline-Parallel Setup:} The CNN runs locally, transmitting activations to Machine~1 and receiving gradients. With 10 update epochs per batch and 32 minibatches per epoch (minibatch size 128), we transmit 4096-dimensional activation and gradient vectors. Total data transfer:
\begin{equation}
100 \times 10 \times 32 \times 128 \times 4096 \times 2 = 32.7\text{B tensors} \approx 66\text{ GB}
\end{equation}

The factor of 2 accounts for bidirectional communication (forward activations and backward gradients). While cloud setup requires 28~GB, pipeline parallelism requires 66~GB, a 2.4$\times$ increase. Table~\ref{tab:data_transfer} summarizes this comparison.

\begin{table}[htbp]
\caption{Communication Overhead Comparison}
\begin{center}
\begin{tabular}{|l|c|c|}
\hline
\textbf{Setup} & \textbf{Data Transferred} & \textbf{Relative Cost} \\
\hline
Cloud Setup & 28 GB & 1.0$\times$ \\
Pipeline-Parallel & 66 GB & 2.4$\times$ \\
\hline
\end{tabular}
\label{tab:data_transfer}
\end{center}
\end{table}

This overhead stems from repeated gradient exchanges during multiple training epochs. In the cloud setup, the entire neural network resides on a single machine, so gradient computation and backpropagation occur entirely in local memory. No gradients are ever transmitted over the network. The only communication cost is the one-time transfer of observations per iteration. Pipeline parallelism, by contrast, splits the network across machines and must exchange activations and gradients for every minibatch across all update epochs, amplifying communication by a factor proportional to the number of training passes over the data.

Our goal is to reduce backward gradient communication while maintaining training performance, making pipeline parallelism viable for bandwidth-constrained environments.
